\subsection{The rotating and Cox pair for the additive pages}
\begin{statement}{P3K.30}{Refuted; physical-line correction}\par\noindent\textbf{Physical-line correction, 17 November 2026.} The displayed central coordinate is refuted. With \(U=M^{-1}\), the retained pair is \((k_U,\ell_U)=(-3H,-24\chi)\), the physical pair is \((3H,24\chi)\), and the physical central class is \(H/2-\chi\). The native free pair \((6\chi,3H-6\chi)\) is retained. K68--69 prove the exact line and transfer maps. The following earlier assertion and its argument are retained as the located historical claim; their corrected current consequences are proved in K68--69.\medskip\par 

For every actual page of P3K.26 and P3K.28, let \(M\) be its page-cover degree,
\(h(G)=1/d\), and \(aG\) its marked line point. Write
\(H=Ma^2/d\) and \(\chi=\chi(\mathcal O_{S_{\rm page}})\).
Use the actual vertical-edge Cox normal product transported by the
extending signed cusp origin. For one positive loop in the actual
parameter \(r\), with \(\lambda=r^M\), the individual Jin pair and
the K23 relative coordinates are
\[
 \begin{aligned}
 k&=-3H,&\ell&=-24\chi,\\
 m_{\rm rel}&=6\chi,&n_{\rm rel}&=3H-6\chi,\\
 s&=\chi-H/2\pmod{28}.
 \end{aligned}
\]
This holds for both one-cap norms \(n=\pm1\) and both joint determinants
\(D=\pm1\). The comparison fixes the actual initial Cox quaternion
frame; a possible extending constant rotation loop changes neither
integer. No different fan edge's relative four-cycle is assigned here.

\end{statement}
\begin{proof}
\par

\medskip\noindent\textbf{1. The two actual cusp products for the signed parameter}

\par\nobreak

Use the physical coordinates of K26 and K28:
\[
 V=e^{2\pi i\zeta_2},\qquad
 B_\lambda(Z,V)=(U(q)^{-1}Z,q\lambda W(q)^{-1}V),
\]
where \(U=e^{2\pi i g}\), \(W=e^{-2\pi i C}\) are the actual
nonvanishing cusp units. Thus the positive outward period is
\[
 \gamma:(Z,V)\longmapsto
      (U(q)Z,q^{-1}\lambda^{-1}W(q)V).
 \tag{K30.2}
\]
At fixed \(|r|\), variation of the phase of \(\lambda=r^M\)
is a physical delta rotation. It disappears in the circle orbit,
including in the two actual unimodular Cox charts. Consequently
the orbit body and its marked period-cell quotient are fixed
before applying the ordinary closing map.

Set \(s=\log q/(2\pi i)\) and \(v=\log r/(2\pi i)\).
The actual extending cusp origin is
\[
 \sigma_\varepsilon(v)
       =\varepsilon\left(gv,-Cv-\frac{M(v^2-v)}2\right).
 \tag{K30.3}
\]
The ordinary parameter deck is \(+\varepsilon s\delta\).
The identity
\[
 \sigma_\varepsilon(v+1)-\sigma_\varepsilon(v)
            =\varepsilon s\delta+\varepsilon\gamma
 \tag{K30.4}
\]
shows that the ordinary deck minus the cusp-origin difference
is the actual integral period \(-\varepsilon\gamma\).
This closes the cusp product, with the same positive \(r\)-loop
as in the construction proofs.

For a fixed initial parameter \(\lambda_0\), the bounded
coordinate map on its universal angular interval is literally
\[
 C_\theta=
 \left(
 U^{\varepsilon\theta},
 (W\lambda_0^{-1})^{\varepsilon\theta}
       e^{-\pi i\varepsilon M\theta(\theta-1)}
 \right).
 \tag{K30.5}
\]
All powers use the specified logarithms. These nonzero bounded
torus maps extend in the actual two adjacent cusp charts and
their Cox master coordinates. At \(\theta=1\) they are exactly
the bounded period-return map from (K30.4). The quadratic
factor is physical delta; it acts trivially on the orbit normal
three-plane. An additional unit in the global ordinary closing
map adds the same bounded origin term to both products and
is carried by the same extending torus map.

The actual ordinary origin is
\[
 o_\varepsilon(s,v)
       =\varepsilon\left(
          -\frac{gs}{M},
          \frac{Cs}{M}+\frac{s^2}{2M}+sv\right).
 \tag{K30.6}
\]
The source and parameter descents of these origins, including
the auxiliary \(2M\)-root gauge, were proved in K26 and K28.
Their literal transition has angular difference
\[
 (\sigma_\varepsilon-o_\varepsilon)(v+\theta)
       -(\sigma_\varepsilon-o_\varepsilon)(v)
   =\varepsilon\theta\gamma
       -\frac{\varepsilon M\theta(\theta-1)}2\,\delta .
 \tag{K30.7}
\]
On the circle orbit this is the marked rotation
\(R_{\varepsilon\theta}\) of the gamma product. Equivalently,
it is the positive rotation \(R^\beta_\theta\) for
\(\beta=\varepsilon\gamma\).

\par

\medskip\noindent\textbf{2. The exact smooth extension and Cox normal comparison}

\par\nobreak

We retain the actual map, rather than differentiate the annular
expression involving \(|q|^{\varepsilon\theta}\) at the node.
In the first small-resolution chart
\[
 (x,v,z),\qquad u=xz,\quad y=vz,
\]
the circle weights on \((v,z)\) are \((-1,+1)\), and its
Hopf coordinates are
\[
                (y,b),\qquad b=(|z|^2-|v|^2)/2.
\]
In the second chart \((u,y,w)\), with \(x=uw,v=yw\), use
\[
                (x,b'),\qquad b'=(|u|^2-|w|^2)/2.
\]
On the overlap, putting \(r=|z|^2\), the exact derivatives are
\[
 2b=r-|y|^2/r,\qquad
 2b'=|x|^2r-r^{-1},\qquad
 \frac{db'}{db}
   =\frac{|x|^2+r^{-2}}{1+|y|^2r^{-2}}>0.
 \tag{K30.8}
\]
The proper increasing orbit coordinate \(\chi_0(x,y,b)\)
of P3F.9, fixed relative to its marked outside collar,
has \((\chi_0)_b>0\). It gives the smooth product map
\(\Phi(x,y,b)=(x,y,\chi_0(x,y,b))\), with the corresponding
second-chart formula.

In the two actual extending products the orbit comparison is
\[
                       J_\theta=R_{\varepsilon\theta}\circ\Phi.
 \tag{K30.9}
\]
Equation (K30.7) proves agreement on the original annular
identification; (K30.8) proves the smooth extension through the
two node charts. The gamma rotation is an actual flow on the
fixed source-preserving product, so its negative direction for
\(\varepsilon=-1\) is equally defined.

Quotienting the derivative by the tangent of the fixed curve,
including the derivative of its graph gamma coordinate, gives
\[
 Q_x=
 \begin{pmatrix}
 1&0&0\\0&1&0\\
 (\chi_0)_{\Re y}&(\chi_0)_{\Im y}&(\chi_0)_b
 \end{pmatrix},
 \qquad(\chi_0)_b>0 .
 \tag{K30.10}
\]
The second chart has the same formula with \(x,b'\).
These maps are independent of \(\theta\). Their source shears
contract linearly, and the positive last entry contracts to one.
The intrinsic normal line is the free gamma direction; the
quotient is the source normal two-plane. Its splittings form
a contractible space, so the displayed contractions agree
globally. This remains true for either signed direction of
the transported rotation in (K30.9).

The sign is also retained explicitly in beta coordinates.
There the product map is
\(\Phi_\beta(x,y,b)=(x,y,\varepsilon\chi_0(x,y,b))\);
its raw normal derivative is
\[
 Q_x^\beta=T_\varepsilon Q_x,\qquad
 T_\varepsilon=\operatorname{diag}(1,1,\varepsilon).
 \tag{K30.10a}
\]
For the negative sign its last raw derivative is indeed negative.
Apply this same constant coordinate and orientation change to
both initial normal frames, as prescribed by the positive-Hopf
Thom orientation in K26 and K28. In the transported reference,
\(T_\varepsilon^{-1}Q_x^\beta=Q_x\), so the remaining comparison
is the positive triangular, theta-independent map just displayed.
Equivalently the full relative clutch is conjugated by
\(T_\varepsilon\); this is also conjugation by
\(\varepsilon T_\varepsilon\in SO(3)\).
This gives the actual placement of the signed coordinate change
at both ends. It contributes no free three-dimensional
framing class.

The initial actual Cox frame is the quaternion frame of K27,
which we retain explicitly. If the fixed-sphere charts satisfy
\(\zeta'=\zeta^{-1}\), \(c_i'=\zeta c_i\), then
\[
 q_{\rm H}
  =(1+j\zeta)(c_0+j\overline{c_1})
  =(\zeta^{-1}+j)(c_0'+j\overline{c_1'}).
 \tag{K30.11}
\]
The two products are literally equal, using \(zj=j\bar z\).
The physical delta action is right multiplication by
\(e^{-i\theta}\), so the positive Hopf normal coordinate is
\(q_{\rm H}i\overline{q_{\rm H}}\), with exactly the K23
orientation. Transport this initial frame by \(J_\theta\).
By (K30.9), it is the same frame obtained by applying
\(dR_{\varepsilon\theta}d\Phi\) to the initial Cox frame.
Equation (K30.10) supplies no parameter-dependent free
normal-framing class.

The bounded units in (K30.5) also retain their actual normal
map. In a Cox chart \(q=\zeta c_0c_1\), a unitary torus map
has form
\[
 (\zeta,c_0,c_1)\longmapsto
       (a\zeta,p c_0,q_1c_1),\qquad apq_1=1.
\]
Choose continuous square roots along the actual interval,
\[
             A^2=pq_1,\qquad B^2=p/q_1,\qquad AB=p.
\]
Substitution in (K30.11) gives the literal equality
\[
                         q_{\rm H}\longmapsto A q_{\rm H}B .
 \tag{K30.12}
\]
Its Hopf normal map is the constant rotation
\(\eta\mapsto A\eta\bar A\), independent of the fixed-sphere
point. The positive factors are the same normal isomorphism
on both sides of the comparison and cancel. The higher
base-unit terms have zero normal linear part, because \(q\)
is quadratic in \(c_0,c_1\). Thus any remaining loop is a
constant rotation loop. It extends over the fixed three-disc
and contributes neither \(k\) nor \(\ell\). There is no
free \(\pi_3(SO(3))\) discrepancy between the actual closed
Cox product and the transported rotating product.

\par

\medskip\noindent\textbf{3. The degree-\(M\) binding coordinate is an actual holomorphic chart}

\par\nobreak

For one cap the map in K26 has weighted polynomial
\(z_2-\tau z_1^M\), and its source-base coordinate is
\(t=z_2/z_1^M\). Its page ray in the binding coordinate
\(q_b=t-\tau\) has exactly
\[
                            \alpha=\phi-M\beta,
 \tag{K30.13}
\]
where \(\phi\) is the fixed outward page phase in radians.

For the two-cap page of K28, its exact ordinary phases give
\[
 \arg(z_2^{m_0}-\tau z_1^{m_\infty})
       =M\beta+\arg(t-\tau)-\varepsilon m_\infty n_0\arg t .
 \tag{K30.14}
\]
Near the regular binding \(t=\tau\ne0\), define the literal
holomorphic coordinate
\[
              q_b=(t-\tau)(t/\tau)^{-\varepsilon m_\infty n_0}.
 \tag{K30.15}
\]
Its derivative at \(\tau\) is one, so it is a local
biholomorphism. The exponent is an integer and the factor is
a specified nonvanishing holomorphic unit. Equations
(K30.14)--(K30.15), after retaining their constant \(\tau\)
phase in \(\phi\), prove exactly (K30.13).
This retains the actual two-cap source coordinate; it does
not replace a variable base phase by its limiting value.

The actual graph of K26 or K28 has
\[
                         \beta=\theta+g(p)
 \tag{K30.16}
\]
on the binding, with \(p\in T^2\) the source point and
\(\theta\) the positive physical parameter angle.
Here \(g:T^2\to\mathbb R/2\pi\mathbb Z\) is the graph phase
fixed by the source-preserving orbit product. Local radian
lifts suffice; \(R_z(-Mg(p))\) is a global map because \(M\)
is an integer. No source-dependent modification of the
page phase is made.

\par

\medskip\noindent\textbf{4. The full normal clutch and its Pontryagin number}

\par\nobreak

Give the graph its oriented page-normal vector and its two
Cartesian parameter vectors. The actual graph derivative
contributes only a contractible triangular shear to this
frame. On the compression choose a normal frame at parameter
zero and transport it by \(R_{\varepsilon\theta}\).
The compression is the same solid torus minus a ball used
in K6. Its tangent bundle is trivial, and
\(TY\oplus\mathbb R\) is trivial for the actual orbit
\(Y=S^5\). Hence its rank-three oriented normal bundle has
\(w_2=0\), and the three-dimensional obstruction groups
give a frame. At the fixed-sphere end its difference from
the specified positive Hopf frame is a map
\(S^2\to SO(3)\), which extends across the fixed three-disc.
Transport this actual extension by the same signed flow.

After quotienting by the source tangent directions and
the graph tangent, use the four coordinates consisting of
the real and imaginary \(q_b\) directions, the parameter
radius, and the relative angle \(\beta-\theta\).
The outward page collar and inward compression tangent
point in the same positive radial half-space. Their actual
common tangent is
\[
           u=(\cos\alpha,\sin\alpha,\rho,0),\qquad\rho>0.
\]
The literal quotient map is
\[
 (v_q,v_r,v_\eta)\longmapsto
       \left(v_q-\frac{v_r}{\rho}(\cos\alpha,\sin\alpha),
               v_\eta\right).
 \tag{K30.17}
\]
It kills \(u\). Project the Cartesian parameter vectors,
whose radial-relative coefficients are
\((\cos\theta,-\sin\theta)\), \((\sin\theta,\cos\theta)\);
reverse both and apply the positive reference rescalings.
The resulting actual oriented graph normal matrix is
\[
 G(\alpha,\theta)=
 \begin{pmatrix}
 -\sin\alpha&\cos\theta\cos\alpha&\sin\theta\cos\alpha\\
 \cos\alpha&\cos\theta\sin\alpha&\sin\theta\sin\alpha\\
 0&\sin\theta&-\cos\theta
 \end{pmatrix}\in SO(3).
 \tag{K30.18}
\]
The signed Thom calculation of K26 and K28 fixes the normal
orientation here as positive Hopf for both signs.
Equivalently, changing from the gamma relative coordinate
to the beta relative coordinate multiplies the reference
three-space by the constant matrix
\(T_\varepsilon=\operatorname{diag}(1,1,\varepsilon)\).
It conjugates both frame comparisons. This conjugation is
also conjugation by the oriented constant matrix
\(\varepsilon T_\varepsilon\in SO(3)\), so it introduces no
additional clutch or Pontryagin term.

The derivatives of \(g(p)\) are exactly the source tangent
shears already removed in (K30.17). The positive \(\rho\)
rescaling is contractible, and the actual collar rounding
stays in the same contractible positive radial half-space.
These specify a homotopy of the actual transition, including
all source tangent terms.

In the transported product coordinates the compression
frame is \(h(p)\), independent of \(\theta\). Since
\(G(\alpha,\theta)=R_z(\alpha)G(0,\theta)\), the full
clutch on the actual \(T^3\) is homotopic to
\[
 \begin{split}
 &h(p)^{-1}G(\phi-M(\theta+g(p)),\theta)\\
 &\quad=
  [h(p)^{-1}R_z(-Mg(p))]
  [R_z(\phi-M\theta)G(0,\theta)] .
 \end{split}
 \tag{K30.19}
\]
The first factor is pulled back from \(T^2\), the second
from the actual parameter circle. Both are globally
defined, including any nonzero one-dimensional frame
classes or graph-phase winding. Expanding
\[
 (ab)^{-1}d(ab)=b^{-1}(a^{-1}da)b+b^{-1}db
\]
and using the Maurer--Cartan identities shows that the
cubic trace form of the product differs from the sum of
the two cubic forms by an exact form. Each factor cubic
form vanishes by dimension. Therefore its integral over
the marked \(T^3\) is zero. The interpolating connection
\(t\,g^{-1}dg\) on the actual collar identifies this integral
with the normal Pontryagin clutching number, and gives
\[
                    \langle p_1(\bar\nu),[\bar X]\rangle=0.
 \tag{K30.20}
\]
This retains the full normal bundle and all lower-degree
frame data; it does not assert that the bundle is trivial.

\par

\medskip\noindent\textbf{5. The individual integers and the relative coordinates}

\par\nobreak

The already proved actual page-line comparison gives
\[
 b^2=4c_1(N)^2=-4H,\qquad
 p_1(T\bar X)=-24\chi .
 \tag{K30.21}
\]
The signed trace bordism has positive meridian integral
for each norm sign. Its sphere-bundle characteristic
formulas, in the matched normal framing just constructed,
are
\[
 k=\frac{p_1(\bar\nu)+3b^2}{4},\qquad
 \ell=p_1(T\bar X)+p_1(\bar\nu).
 \tag{K30.22}
\]
Substitution of (K30.20)--(K30.21) proves
\(k=-3H\), \(\ell=-24\chi\).
The actual line and height equality
\(H=2\chi+2(P\cdot O)\) makes both integers integral.
The full K23 coordinate morphism is
\[
 m_{\rm rel}=-k+6A=-\ell/4,\qquad
 n_{\rm rel}=-6A,\qquad s=A\bmod28 .
 \tag{K30.23}
\]
Together with the independently proved
\(A=\chi-H/2\), this gives (K30.1).

The construction fixes the initial Cox quaternion product,
the signed cusp product, and its actual transport. The
possible constant rotation loop extends over the fixed
three-disc and changes neither number. The formula is for
one positive loop in the stated \(r\) parameter; it has not
been divided by the page degree, by a source-root order,
or by the number of scalar turns in \(\lambda\).

\par

\medskip\noindent\textbf{Certificate}

\par\nobreak

The exact source-phase and normal-matrix identities are
executed in \path{code/part3/general_cap_framing.py}. It also
combines the already certified K26 and K28 page arithmetic
with (K30.22)--(K30.23). The actual normal comparison is
the explicit map and clutch proof above.

\end{proof}
